
\documentclass[../../main.tex]{subfiles}
\begin{document}

% 年度の大きな表示
\begin{center}
  {\fontsize{48pt}{58pt}\selectfont \textbf{1989}\normalsize\textbf{年度}}
\end{center}

\vspace{10mm}

% 本文

\vspace{15mm}

% 出題テーマと難易度の表
\noindent\textbf{出題テーマと難易度}

\vspace{5mm}

\begin{center}
  \shadowbox{%
    \begin{tabular}{|c|p{8cm}|c|}
      \hline
      \textbf{問題番号} & \textbf{テーマ} & \textbf{難易度} \\
      \hline
      第1問 & 放物線の接線と面積 & \\
      \hline
      第2問 & エピサイクロイド（曲線の長さ） & \\
      \hline
      第3問 & 関数方程式と微分可能性 & \\
      \hline
      第4問 & 定積分の極限 & \\
      \hline
      第5問 & 確率と期待値（漸化式） & \\
      \hline
    \end{tabular}%
  }
\end{center}

\vspace{15mm}

% 解答
\noindent\textbf{解答}

\vspace{5mm}

\begin{center}
  \shadowbox{%
    \renewcommand{\arraystretch}{1.6}
    \begin{tabular}{|c|c|p{8cm}|}
      \hline
      \textbf{問題番号} & \textbf{小問} & \textbf{答え} \\
      \hline
      第1問 & (2) & $Y=X^2+1$ \\
      \hline
      \multirow{2}{*}{第2問} & (1) & $(2\sqrt{2},\pm\sqrt{2})$ \\
      \cline{2-3}
                             & (2) & $L=24$ \\
      \hline
      \multirow{2}{*}{第3問} & (1) & 証明問題 \\
      \cline{2-3}
                             & (2) & $f(x)=0$ または $f(x)=e^{f'(0)x}-1$ \\
      \hline
      第4問 & - & $\dfrac{2}{3}\pi^2$ \\
      \hline
      第5問 & (2) & $E_k=n\left[1-\left(\dfrac{n-1}{n}\right)^k\right]$ \\
      \hline
    \end{tabular}%
    \renewcommand{\arraystretch}{1.0}
  }
\end{center}

\vspace{15mm}

% 解答時間
\noindent\textbf{試験形態}

\vspace{5mm}

\begin{center}
  \shadowbox{%
    \begin{tabular}{p{8cm}}
      （要確認）
    \end{tabular}%
  }
\end{center}

\end{document}
